\section{All-time momentum kernels of the true marked clearance source}
\label{sec:nm-kernel}

The first-image source theorem and its admissible growth theorem are
already established in Sections~\ref{sec:ag-source} and
\ref{sec:uc-growth}.  We retain those proofs and constants.  The new
conclusions below concern every later or earlier physical momentum,
joint clearance--incidence control, and the one-sided marked source at
zero clearance.  They do not identify momentum with a Markov state.

All radii belong to $[9/20,47/100]$, the speed is one, and
$\delta_*=3/200$.  Write $a_\circ=\sqrt{1827}/200$ for the
foot standoff of Lemma~\ref{lem:xf-standoff}, $k_-=100/47$ and
$t_*=3/50$.  The constants $M_*,A_*,C_g,J_f$ have exactly their
meanings in Sections~\ref{sec:ag-source}--\ref{sec:uc-growth};
$A_*$ includes both impact signs.  Fix
\begin{equation}\label{eq:nm-H}
 H_{\rm nm}=\max\left\{\frac{2M_*}{k_-a_\circ},
                    \frac{A_*C_gJ_f}{k_-t_*}\right\}.
\end{equation}
No curve-class enlargement is needed: these are the original
angular-source graphs, not the different initial momentum fibres of
Section~\ref{sec:mc-fibres}.

For $u>0$, start at the actual marked foot
$-\sigma(R+u)v_\vartheta^\perp$, with direction $v_\vartheta$,
and let $P_l^+(\vartheta,\sigma,u)$ be the tangential momentum at
the $l$th actual collision after that foot, $l\ge1$.  Let $P_l^-$
be the tangential momentum in the original forward record at the
$l$th collision before the foot.  Thus its sign is opposite to the
outgoing momentum of the time-reversed record.  All visited disks
and all intervening flights are the actual ones.  There is no
collision, reflection, or change of the physical clock at the mark.

\begin{theorem}[A uniform angular kernel at every actual collision]
\label{thm:nm-kernel}
For every $R$, $0<u\le\delta_*$, $l\ge1$, Borel set
$I\subset(-1,1)$ and source selectors $0\le h_\sigma\le1$,
\begin{equation}\label{eq:nm-kernel}
 \sum_{\sigma=\pm1}\int_0^{2\pi}
    h_\sigma(\vartheta)\one_{\{P_l^\pm\in I\}}\,d\vartheta
                   \le H_{\rm nm}|I|.
\end{equation}
The integrals are over the regular finite records; their complement
is angular-null.  For either direction and every $v\in(-1,1)$,
the complete sum over distinct regular angular roots satisfies
\begin{equation}\label{eq:nm-roots}
 \sum_{\sigma=\pm1}
 \sum_{\vartheta:\,P_l^\pm(\vartheta,\sigma,u)=v}
       \frac{1}{|\partial_\vartheta P_l^\pm|}
                    \le H_{\rm nm}.
\end{equation}
Artificial seams and subdivisions do not create extra roots.
Selectors may retain any of the original exact records or original
physical-roof Borel tests.  The constants are independent of $l$,
the full collision words, the radius, clearance, and selectors.
\end{theorem}
\begin{proof}
For the first true collision, Theorem~\ref{thm:ag-density} gives
$ds/d\vartheta=-a/c$ and $d\varphi/ds=\kappa+c/a$.
Consequently
\[
 \frac{dP_1^+}{d\vartheta}
       =c\frac{d\varphi}{ds}\frac{ds}{d\vartheta}
       =-(\kappa a+c),\qquad
 |\partial_\vartheta P_1^+|\ge k_-a_\circ.
\]
There are at most $M_*$ first-hit intervals for each sign.  On every
such interval the momentum is strictly monotone.  Direct substitution
therefore proves \eqref{eq:nm-kernel} for $l=1$ with the first
constant in \eqref{eq:nm-H}, without cutting the target momentum
into infinitely many grazing strips.

For $l\ge2$, use the exact, source-labelled first-image measure
of Theorem~\ref{thm:ag-family}.  Its sum of density suprema is
$A_*$.  Propagate each of these homogeneous short unstable curves
through exactly $l-2$ further physical collisions.  Theorem~
\ref{thm:uc-growth} bounds the complete sum of inverse arclength
Jacobians by $C_g$.  On each resulting incoming short graph, leave
one actual collision to perform the terminal coarea.  The exact
cancellation \eqref{eq:mc-last-cancellation} gives
\[
 \left|\frac{dp_{\rm next}}{d\ell}\right|\ge k_-t_*,
 \qquad
 \int_C\one_{\{p_{\rm next}\in I\}}\,d\ell
                   \le\frac{J_f}{k_-t_*}|I|.
\]
The small terminal incidence cancels before estimating; the last
coarea is cut into its finitely many physical next-hit branches,
with no terminal homogeneous subdivision.  Sum the labelled
densities and all intermediate growth pieces by Tonelli.  This gives
$A_*C_gJ_f/(k_-t_*)$ and proves the forward inequality.  A Borel
selector is retained until the positive upper bound and incurs no
component count or derivative.

At each finite stage the singular-null coverage of
Theorem~\ref{thm:uc-growth}, followed by the finite last-hit cuts,
applies to the absolutely continuous labelled source.  The omitted
angular set is null.  The momentum derivative is nonzero on every
regular complete branch.  Coarea gives the root formula almost
everywhere in $v$.  To obtain the bound at every regular root value,
fix a finite set of distinct roots at $v_0$, choose disjoint regular
source intervals, and use their continuous positive inverse
derivatives.  Their sum is bounded almost everywhere in a common
neighborhood of $v_0$, hence also at $v_0$.  Taking the supremum
over finite root sets proves \eqref{eq:nm-roots}.

Finally time reversal sends
$(\vartheta,\sigma,u)$ to $(\vartheta+\pi,-\sigma,u)$ and
preserves $d\vartheta$ and the two-sign sum.  On collision
coordinates it is $\mathcal I(\alpha,p)=(\alpha,-p)$.
Thus $P_l^-=-P_l^+(\vartheta+\pi,-\sigma,u)$ in the reversed
record.  Apply the forward result to $-I$, carrying the original
selector through this bijection.  This does not replace the
original section by its reversal: the section tests are still
inside that transported selector.
\end{proof}

For a Borel clearance set $E\subset(0,\delta_*]$, let $N_E(x)$
count the distinct nonincident disks whose original flight has
clearance in $E$.  An interval convention at its endpoints is
immaterial to the collision-area identity, but will be retained in
the one-sided trace below.  For an $m$-flight record and its marked
flight $j$, write $x=X_{R,j}=T_R^{-j}X_R$ as in
Section~\ref{sec:cf-record}.  At collision index $0\le i\le m$,
the momentum is obtained after exactly $j-i+1$ backward collisions
when $i\le j$, and after $i-j$ forward collisions when $i\ge j+1$.
These are both positive integers, also at the two record endpoints.

\begin{corollary}[Joint clearance and any physical collision momentum]
\label{cor:nm-joint}
For every $m\ge1$, $0\le j<m$, $0\le i\le m$, every original
exact event $\Pi=\Pi_{n,k,m,R}$, $0\le h\le1$, and Borel sets
$E\subset(0,\delta_*]$, $I\subset(-1,1)$,
\begin{equation}\label{eq:nm-joint}
 \int h\one_\Pi N_E(T_R^jx)\one_{\{p_i\in I\}}\,d\nu_R^*(x)
             \le\frac{H_{\rm nm}}{4\pi R c_*}|E|\,|I|.
\end{equation}
The same domination holds with any bounded Borel original-roof
restriction, including an arbitrarily small physical time interval.
It is a joint domination of a marked counting measure, not a
conditional independence statement.
\end{corollary}
\begin{proof}
Use invariance of the original collision area to apply the exact
one-copy flux identity \eqref{eq:cf-flux} at the marked flight
$j$.  Transport $h\one_\Pi$ and all its true record information
to $X_{R,j}$.  Conditioning contributes exactly $c_*^{-1}$, and
the physical line Jacobian contributes exactly $(4\pi R)^{-1}$.
For each fixed $u\in E$, use Theorem~\ref{thm:nm-kernel} with
the positive collision count just specified for $p_i$.  Integrate
the resulting bound in $u$.  Tonelli first on simple clearance
sets and then on their Borel extension proves the statement.
Every distinct marked disk is counted once by the flux bijection;
no collision-word multiplicity is estimated by hand.
\end{proof}

\begin{corollary}[Uniform singular-incidence moments on every clearance source]
\label{cor:nm-moments}
Let
\[
 w_i=q_0^{\min(i,m-i)},\quad q_0=47/53,\quad
 I_m=1+\sum_{i=1}^{m-1} w_i/c_i,\quad
 W_0=1+\frac{2}{1-q_0}.
\]
For $1\le\gamma<2$, define the finite number
\[
 J_\gamma=\int_{-1}^1(1-v^2)^{-\gamma/2}\,dv,
 \qquad
 B_\gamma=4\pi+\frac{2H_{\rm nm}J_\gamma}{1-q_0}.
\]
Then, uniformly in every original exact label, $m,j,R$, and
$0<\delta\le\delta_*$,
\begin{align}
 \int h\one_\Pi N_\delta(T_R^jx)I_m^\gamma\,d\nu_R^*
  &\le\frac{\delta W_0^{\gamma-1}B_\gamma}{4\pi R c_*},
                                             \label{eq:nm-moment}\\
 \int h\one_\Pi N_\delta(T_R^jx)
                         \one_{\{c_i<\eta\}}\,d\nu_R^*
  &\le\frac{H_{\rm nm}\delta}{2\pi R c_*}\eta^2,
       \qquad 0<\eta\le1.                  \label{eq:nm-grazing}
\end{align}
For arbitrary nonnegative depth coefficients $b_j$ and widths
$0<\delta_j\le\delta_*$, the right side of
\eqref{eq:nm-moment}, summed against $b_j$, is obtained by
replacing $\delta$ by $\sum_j b_j\delta_j$.
This includes every shallow flight, with no incidence-load cutoff.
\end{corollary}
\begin{proof}
For each fixed angular clearance fibre, Theorem~\ref{thm:nm-kernel}
and monotone approximation give
\[
 \sum_\sigma\int c_i^{-\gamma}\,d\vartheta
                        \le H_{\rm nm}J_\gamma.
\]
The integral $J_\gamma$ is finite precisely in the stated range:
near either endpoint $1-v^2$ is comparable to $1-|v|$.
The angular mass of both signs is $4\pi$.  Weighted convexity
with the extra weight one for the constant term gives
\[
 I_m^\gamma\le
 \left(1+\sum_{i=1}^{m-1}w_i\right)^{\gamma-1}
                 \left(1+\sum_{i=1}^{m-1}w_ic_i^{-\gamma}\right).
\]
Since $\sum_iw_i\le2/(1-q_0)$, integrating this bound on the
complete angular source gives $W_0^{\gamma-1}B_\gamma$.
Apply the exact flux and its original conditional factor.  This
proves \eqref{eq:nm-moment}.  The set $c_i<\eta$ has momentum
length $2(1-\sqrt{1-\eta^2})\le2\eta^2$; insert this set into
\eqref{eq:nm-joint} to prove \eqref{eq:nm-grazing}.
The last statement follows by summing nonnegative measures; it
requires neither independence nor separation of the marked flights.
\end{proof}

\subsection{A genuine one-sided unvisited-tangency source}

The next theorem defines a boundary measure, not a new value of the
collision map at a singular point.  Fix $R,m,j$.  Let $u\downarrow0$
in the positive-clearance family above.  A limiting strict record
lists the actual visited collisions of this family.  The marked
tangency at the foot stays an unvisited boundary event.  At $u=0$
we never use the putative collision map obtained by adding or
deleting a reflection at that foot.

\begin{theorem}[Exact marked trace with original finite record tests]
\label{thm:nm-trace}
For angular-almost every $(\vartheta,\sigma)$ there is an open
one-sided neighborhood in $u$ on which the fixed finite forward
and backward strict word is constant and all its visited collisions
are regular.  Its contacts, momenta and physical roof extend
real-analytically to $u=0$.  Every original finite record selector
$\one_{\Pi_{n,k,m,R}}$ has a one-sided limit there.  The same is
true after multiplying by the indicator of any specified physical
roof interval, with its original open or closed endpoint convention.

Denote the limiting momentum by $p_i^0$ and the one-sided selector
by $\pi^+$.  For any bounded continuous function $f$ of the finite
physical record, including its actual roof, the positive marked
trace is the actual limit
\begin{equation}\label{eq:nm-trace}
 \lim_{u\downarrow0}\frac1{4\pi R c_*}
   \sum_\sigma\int f(\mathrm{record}_u)\pi_u\,d\vartheta
 =\frac1{4\pi R c_*}
   \sum_\sigma\int f(\mathrm{record}_0)\pi^+\,d\vartheta.
\end{equation}
More generally the equality holds for any bounded test with an
angular-almost-everywhere one-sided limit.  The endpoint kernels
\eqref{eq:nm-kernel} and the angular moment bounds used in
\eqref{eq:nm-moment} hold at $u=0$ as well, with the same constants.
In particular the marked trace marginal satisfies
\begin{equation}\label{eq:nm-trace-marginal}
 (p_i^0)_\#\left(\frac{\pi^+}{4\pi R c_*}
                          \sum_\sigma d\vartheta\right)
             \le\frac{H_{\rm nm}}{4\pi R c_*}\,dv.
\end{equation}
No null set common to uncountably many radii is asserted.
\end{theorem}
\begin{proof}
At $u=0$ the foot has distance $R$ from its marked disk and
distance at least $1-R\ge1-R_+>R_+$ from every other center.
The two open tangent rays miss the interior of the marked disk.
The same finite candidate geometry and standoff proof as in
Theorem~\ref{thm:ag-density} applies to their first nonmarked
true hits.  Except for finitely many angular tangencies and seams,
these hits are regular.  Their derivatives, first-image densities,
and curvature bounds are the continuous one-sided versions of
\eqref{eq:ag-density}--\eqref{eq:ag-regularity}; the strict
standoff $a\ge a_\circ$ is unchanged.  The first family is
therefore the same admissible homogeneous class, with the same
summable allowance $A_*$.

Iterate each first-image branch through the finite number of
required true hits in both directions.  At each stage the physical
growth partition removes a countable angular-null set of tangencies
and artificial boundaries, exactly as in the proof of
Theorem~\ref{thm:nm-kernel}.  Away from this set all finitely many
visited incidences are positive.  On each such designated record,
the line--circle intersection with nonzero incidence and the
specular reflection formula are real-analytic functions of the
incoming state and $u$.  After the first true hits, the marked disk is restored as an ordinary
possible next scatterer.  Any later tangency with it or another disk
is part of the same growth singular set.  Every nonincident disk
therefore has positive clearance on this finite strict record, except
for the prescribed original foot tangency;
the finite-horizon candidate reduction makes this a finite
visibility check.  Hence the same true word persists for all
sufficiently small positive $u$.  The marked flight continues to
miss its marked disk because its prescribed clearance is $u>0$.
This establishes the asserted one-sided analytic record.  It uses
no common neighborhood size over words, $m$, or radii.

The lattice displacement is constant on this true word.  Each of
the $m+1$ original section tests is a finite Boolean combination
of strict inequalities in the real-analytic coordinates
$\alpha_i(u),p_i(u)$, using a local circular lift as necessary.
The visit count includes $i=0,\ldots,m-1$ and the terminal test
remains separate.  A real-analytic function of $u$ at zero is
either identically zero near zero or has the sign of its first
nonzero Taylor coefficient for all sufficiently small positive
$u$.  Thus every such strict or nonstrict inequality has a
one-sided indicator limit.  The same argument applies to
$L_m(u)-a$ and $L_m(u)-b$ for any specified roof interval with
endpoints $a,b$.  If a boundary equality persists identically, its
membership is exactly its prescribed endpoint convention.  This
does not identify the trace with evaluation of an arbitrarily
redefined section on a singular record.

All limits in \eqref{eq:nm-trace} now follow by dominated
convergence on the finite angular measure.  To prove the kernel
at $u=0$, one can repeat the already justified first-image,
growth, and last-collision coarea proof on these limiting strict
branches.  Alternatively, apply the $u>0$ domination to
nonnegative continuous momentum tests and pass to zero; finite
measure regularity then gives domination for every Borel set.
Finally use monotone truncation of $c_i^{-\gamma}$ and the same
weighted convexity inequality to obtain the angular moment bound.
All tests with values in $[0,1]$ only restrict a positive measure.
\end{proof}

\paragraph{The new scope and the remaining normalization.}
Theorem~\ref{thm:nm-trace} supplies the complete labelled
one-sided angular source needed for an unvisited clearance
boundary.  Its measure is positive; angular-null exceptional
parameters are not assigned zero charge in an unrelated endpoint
pairing.  The trace is not the original roof--Gauss-angle endpoint
charge.  The exact conversion to that charge includes its physical
normal speed, which is proved in the next section.  Neither the
kernel nor the moment bounds insert an $m^{-2}$ factor into a
physical-roof density.  The original supremum, essential supremum
and ordered limit in \eqref{eq:sc-original-obligation} are retained.
